\documentclass[12pt,a4paper]{article}
\usepackage[utf8]{inputenc}
\usepackage{geometry}
\usepackage{enumitem}
\geometry{margin=2cm}

\title{Questionário – Tema 6: Critérios de Elegibilidade para Revisão Sistemática}
\author{Departamento de Física, Universidade Eduardo Mondlane}
\date{\today}

\begin{document}

\maketitle
\newpage

\section*{Instruções}
\begin{itemize}
    \item Marque a alternativa correta para cada questão.
    \item Algumas questões podem ter mais de uma resposta correta.
    \item Utilize conhecimento adquirido na aula sobre critérios de elegibilidade e PICOTS.
\end{itemize}

\newpage
\section*{Questões}

\begin{enumerate}

\item O que definem os critérios de elegibilidade em uma revisão sistemática?
\begin{enumerate}
    \item A metodologia estatística usada nos estudos
    \item Quais estudos serão incluídos ou excluídos
    \item O número de autores dos estudos
    \item Apenas o idioma dos artigos
\end{enumerate}

\item Qual dos seguintes exemplos é um critério de elegibilidade típico?
\begin{enumerate}
    \item Cor do fundo do gráfico dos resultados
    \item Número de páginas do artigo
    \item Tipo de estudo (ensaios clínicos, observacionais, revisões, etc.)
    \item Nome do orientador da pesquisa
\end{enumerate}

\item Critérios amplos na revisão sistemática tendem a:
\begin{enumerate}
    \item Incluir mais estudos, aumentar a heterogeneidade, reduzir risco de omitir dados
    \item Incluir menos estudos e aumentar a homogeneidade
    \item Tornar os resultados inconsistentes
    \item Excluir apenas ensaios clínicos
\end{enumerate}

\item Critérios restritos na revisão sistemática tendem a:
\begin{enumerate}
    \item Incluir mais estudos, aumentar heterogeneidade
    \item Maior homogeneidade, resultados mais consistentes, risco de perder evidências
    \item Não considerar desfechos relevantes
    \item Selecionar apenas estudos em língua portuguesa
\end{enumerate}

\item Qual destes itens faz parte do PICOTS?
\begin{enumerate}
    \item P – População
    \item I – Intervenção
    \item C – Comparação
    \item O – Desfecho
    \item T – Tempo de seguimento
    \item S – Tipo de estudo
    \item L – Localização geográfica do autor
\end{enumerate}

\item Quais tipos de viés podem ocorrer na seleção de estudos?
\begin{enumerate}
    \item Viés de publicação
    \item Viés de idioma
    \item Viés de tempo
    \item Viés de gráficos coloridos
\end{enumerate}

\item Ao definir critérios de elegibilidade, qual é a vantagem de utilizar critérios amplos?
\begin{enumerate}
    \item Captar maior número de estudos relevantes e reduzir risco de omitir dados
    \item Garantir resultados totalmente homogêneos
    \item Selecionar apenas estudos em inglês
    \item Evitar qualquer heterogeneidade
\end{enumerate}

\item Qual tipo de estudo tem maior nível de evidência?
\begin{enumerate}
    \item Observacional
    \item Ensaios clínicos randomizados (ECR)
    \item Relatórios não-ingleses
    \item Literatura cinzenta
\end{enumerate}

\item Estudos em línguas diferentes do inglês podem ser incluídos se:
\begin{enumerate}
    \item Traduzidos e considerados relevantes
    \item Apenas se forem revisões sistemáticas
    \item Ignorados sempre
    \item Publicados antes de 2000
\end{enumerate}

\item Sobre o refinamento de critérios de elegibilidade, é correto afirmar que:
\begin{enumerate}
    \item Deve ser contínuo durante a revisão para minimizar viés
    \item Só é necessário ao final da revisão
    \item Não é importante em revisões sistemáticas
    \item Depende apenas do número de estudos incluídos
\end{enumerate}

\end{enumerate}

\end{document}
